% Draf terjemahan Bahasa Indonesia dari chapter2.tex baris 1564--1754.
% Sumber: Methods of Algebra, Volume 2, commit
% 9a5803ff77dd3257484cb177f851a73770a59dd3 (CC BY 4.0).
% stable-unit-id: o014.aljabr2.chapter2.serre-subcategories-and-k0-groups
% Segmen ini berhenti tepat sebelum section sumber pada baris 1756.

% segment-id: o014.li2.chapter2.serre-subcategories-and-k0-groups.h001
\section{Subkategori Serre dan Grup \texorpdfstring{$\mathrm{K}_0$}{K0}}\label{sec:Serre-subcat}
\hypertarget{o014.aljabr2.chapter2.serre-subcategories-and-k0-groups}{ }

% segment-id: o014.li2.chapter2.serre-subcategories-and-k0-groups.p001
Subkategori penuh dari kategori abelian secara otomatis mewarisi struktur
kategori-$\cate{Ab}$, sehingga kita dapat menanyakan apakah subkategori penuh
tersebut memiliki sifat kategori aditif atau kategori abelian.

% segment-id: o014.li2.chapter2.serre-subcategories-and-k0-groups.d001
\begin{definisi}\index{kategori abelian!subkategori abelian (abelian subcategory)}
Jika $\mathcal{B}$ merupakan subkategori penuh dari kategori abelian $\mathcal{A}$,
$\mathcal{B}$ sendiri merupakan kategori abelian, dan funktor inklusi
$\iota:\mathcal{B}\to\mathcal{A}$ eksak (Definisi~\ref{def:exact-functor}), maka
$\mathcal{B}$ disebut \emph{subkategori abelian} dari $\mathcal{A}$.
\end{definisi}

% segment-id: o014.li2.chapter2.serre-subcategories-and-k0-groups.pr001
\begin{proposisi}\label{prop:Abel-subcat}
Subkategori penuh $\mathcal{B}$ dari kategori abelian $\mathcal{A}$ merupakan
subkategori abelian jika dan hanya jika kondisi berikut berlaku.
% segment-id: o014.li2.chapter2.serre-subcategories-and-k0-groups.l001
\begin{itemize}
% segment-id: o014.li2.chapter2.serre-subcategories-and-k0-groups.i001
	\item $0\in\Obj(\mathcal{B})$;
% segment-id: o014.li2.chapter2.serre-subcategories-and-k0-groups.i002
	\item jika $X,Y\in\Obj(\mathcal{B})$, maka $X\oplus Y$ juga dapat dipilih
	berada di $\mathcal{B}$;
% segment-id: o014.li2.chapter2.serre-subcategories-and-k0-groups.i003
	\item untuk setiap morfisme $f:X\to Y$, jika $X,Y\in\Obj(\mathcal{B})$, maka
	$\Ker(f)$ dan $\Coker(f)$ juga dapat dipilih berada di $\mathcal{B}$.
\end{itemize}
\end{proposisi}

% segment-id: o014.li2.chapter2.serre-subcategories-and-k0-groups.pf001
\begin{bukti}
Perhatikan bahwa kondisi-kondisi ini bersifat swa-dual. Arah “hanya jika” jelas.
Sekarang andaikan kondisi-kondisi di atas berlaku. Maka $\mathcal{B}$ merupakan
kategori aditif, sedangkan $\iota:\mathcal{B}\to\mathcal{A}$ merupakan funktor
aditif. Selanjutnya kita tunjukkan bahwa $\iota$ mempertahankan semua
$\varprojlim$ hingga dan $\varinjlim$ hingga. Pertama, $\iota$ mempertahankan
$0$ dan jumlah langsung hingga. Selain itu, penyama dapat direpresentasikan
menggunakan $\Ker$ (Catatan~\ref{rem:equalizer-Ker}), sedangkan kondisi tersebut
menyatakan bahwa subkategori $\mathcal{B}$ tertutup terhadap pengambilan $\Ker$.
Jadi $\iota$ mempertahankan semua $\varprojlim$ hingga; kasus $\varinjlim$
bersifat dual.

Dengan mengingat kembali definisi $\Image$ dan $\Coim$ pada
Definisi~\ref{def:Im-Coim}, serta menggunakan langkah sebelumnya, kita tahu bahwa
$\mathcal{B}$ juga tertutup terhadap kedua konstruksi tersebut. Untuk setiap
morfisme $f$ di $\mathcal{B}$, morfisme kanonik
$\Coim(f)\rightiso\Image(f)$ pada diagram \eqref{eqn:strict-morphism} merupakan
isomorfisme di $\mathcal{A}$, sehingga juga merupakan isomorfisme di
$\mathcal{B}$. Jadi $\mathcal{B}$ merupakan subkategori abelian.
\end{bukti}

% segment-id: o014.li2.chapter2.serre-subcategories-and-k0-groups.d002
\begin{definisi}[J.-P.\ Serre]\label{def:Serre-subcat}
\index{subkategori Serre (Serre subcategory)}
\index{subkategori Serre lemah (weak Serre subcategory)}
Subkategori penuh $\mathcal{T}$ dari kategori abelian $\mathcal{A}$ disebut
\emph{subkategori Serre} dari $\mathcal{A}$ jika memenuhi kondisi berikut.
% segment-id: o014.li2.chapter2.serre-subcategories-and-k0-groups.l002
\begin{itemize}
% segment-id: o014.li2.chapter2.serre-subcategories-and-k0-groups.i004
	\item $0\in\Obj(\mathcal{T})$;
% segment-id: o014.li2.chapter2.serre-subcategories-and-k0-groups.i005
	\item untuk setiap barisan eksak pendek
	$0\to X'\to X\to X''\to0$ di $\mathcal{A}$, berlaku
	$X\in\Obj(\mathcal{T})$ jika dan hanya jika
	$X',X''\in\Obj(\mathcal{T})$.
\end{itemize}
Jika kondisi terakhir dilonggarkan menjadi: untuk setiap barisan eksak
di $\mathcal{A}$
% segment-id: o014.li2.chapter2.serre-subcategories-and-k0-groups.q001
\[W\to X'\to X\to X''\to Y,\]
jika $W,X',X'',Y\in\Obj(\mathcal{T})$, maka $X\in\Obj(\mathcal{T})$,
maka $\mathcal{T}$ disebut \emph{subkategori Serre lemah} dari $\mathcal{A}$
\footnote{Definisi yang sedikit mendadak ini dirancang khusus untuk kategori
turunan; lihat \S\sourcecrossref{sec:derived-cat}{4.4}.}.
\end{definisi}

% segment-id: o014.li2.chapter2.serre-subcategories-and-k0-groups.p002
Sebagai contoh, untuk gelanggang komutatif $R$, semua modul Noetherian
(atau Artinian) membentuk subkategori Serre dari $R\dcate{Mod}$.
Ini adalah fakta standar dalam teori modul \cite[Lema 6.10.2]{Li1}.

% segment-id: o014.li2.chapter2.serre-subcategories-and-k0-groups.p003
Jika $\mathcal{T}$ merupakan subkategori Serre (atau subkategori Serre lemah)
dari $\mathcal{A}$, maka $\mathcal{T}^{\opp}$ merupakan subkategori Serre
(atau subkategori Serre lemah) dari $\mathcal{A}^{\opp}$.

% segment-id: o014.li2.chapter2.serre-subcategories-and-k0-groups.p004
Subkategori Serre lemah $\mathcal{T}$ memiliki sifat jenuh berikut: jika
$X\in\Obj(\mathcal{A})$ isomorfik dengan suatu objek dari $\mathcal{T}$, maka
$X\in\Obj(\mathcal{T})$; ini merupakan konsekuensi menerapkan kondisi terakhir
pada barisan eksak $0\to0\to X\rightiso Y\to0$.

% segment-id: o014.li2.chapter2.serre-subcategories-and-k0-groups.co001
\begin{korolari}\label{prop:Serre-subcat-abelian}
Jika $\mathcal{T}$ merupakan subkategori Serre lemah dari kategori abelian
$\mathcal{A}$, maka $\mathcal{T}$ merupakan subkategori abelian dari
$\mathcal{A}$.
\end{korolari}

% segment-id: o014.li2.chapter2.serre-subcategories-and-k0-groups.pf002
\begin{bukti}
Cukup memverifikasi kondisi dalam Proposisi~\ref{prop:Abel-subcat}. Pertama,
$0\in\Obj(\mathcal{T})$. Selanjutnya, masukkan barisan eksak berikut ke dalam
definisi subkategori Serre lemah:
% segment-id: o014.li2.chapter2.serre-subcategories-and-k0-groups.g001
\begin{gather*}
	0\to X\to X\oplus Y\to Y\to0
	\quad\text{(Proposisi~\ref{prop:biproduct-ses})},\\
	0\to0\to\Ker(f)\to X\xrightarrow{f}Y,\quad
	X\xrightarrow{f}Y\to\Coker(f)\to0\to0,
\end{gather*}
sehingga terlihat bahwa $\mathcal{T}$ tertutup terhadap jumlah langsung,
$\Ker$, dan $\Coker$.
\end{bukti}

% segment-id: o014.li2.chapter2.serre-subcategories-and-k0-groups.ex001
\begin{contoh}\index[sym1]{kerF@$\Ker(F)$}
Jika $F:\mathcal{A}\to\mathcal{B}$ merupakan funktor eksak antara kategori
abelian, maka semua objek yang memenuhi $FX=0$ membentuk subkategori Serre
dari $\mathcal{A}$, yang dinotasikan dengan $\Ker(F)$.
\end{contoh}

% segment-id: o014.li2.chapter2.serre-subcategories-and-k0-groups.th001
\begin{teorema}[Kuosien Serre]\label{prop:Serre-quotient}
\index{kuosien Serre (Serre quotient)}
Misalkan $\mathcal{T}$ merupakan subkategori Serre dari kategori abelian
$\mathcal{A}$. Maka terdapat kategori abelian $\mathcal{A}/\mathcal{T}$
(yang mungkin merupakan kategori besar), bersama funktor eksak yang pada
hakikatnya surjektif $Q:\mathcal{A}\to\mathcal{A}/\mathcal{T}$, sedemikian
sehingga $\Ker(Q)=\mathcal{T}$, dan berlaku sifat universal berikut: untuk
setiap kategori abelian $\mathcal{B}$ dan setiap funktor eksak
$F:\mathcal{A}\to\mathcal{B}$ dengan $\Ker(F)\supset\mathcal{T}$, terdapat
tepat satu funktor eksak $G:\mathcal{A}/\mathcal{T}\to\mathcal{B}$
sedemikian sehingga $F=GQ$.

Dalam sifat universal ini, $G$ merupakan funktor setia jika dan hanya jika
$\Ker(F)=\mathcal{T}$.
\end{teorema}

% segment-id: o014.li2.chapter2.serre-subcategories-and-k0-groups.pf003
\begin{bukti}
Definisikan
$S:=\left\{f\in\Mor(\mathcal{A}):\Ker(f),\Coker(f)\in\Obj(\mathcal{T})\right\}$.
Kita klaim bahwa $S$ merupakan sistem multiplikatif dalam pengertian
Definisi~\ref{def:multiplicative-family}.

Jelas bahwa $S$ memuat semua morfisme identitas, sehingga (S1) berlaku.
Misalkan $f:X\to Y$ dan $g:Y\to Z$ merupakan elemen $S$. Berdasarkan
barisan eksak yang saling dual,
% segment-id: o014.li2.chapter2.serre-subcategories-and-k0-groups.g002
\begin{gather*}
	0\to\Ker(f)\to\Ker(gf)\xrightarrow{f}\Ker(g),\\
	\Coker(f)\xrightarrow{g}\Coker(gf)\to\Coker(g)\to0,
\end{gather*}
segera diperoleh $gf\in S$, sehingga (S2) berlaku. Selanjutnya pertimbangkan
morfisme $X\xrightarrow{s\in S}Z\xleftarrow{f}Y$. Definisikan
$W:=X\dtimes{Z}Y$, dengan morfisme $s':W\to Y$. Perhatikan bahwa
$\Ker(s')\simeq\Ker(s)$ (Proposisi~\ref{prop:pull-back-ker}), dan $f$
menginduksi $\Coker(s')\hookrightarrow\Coker(s)$
(Korolari~\ref{prop:Abel-cat-Coker-pull}); maka $s'\in S$. Jadi (S3) berlaku.
Akhirnya, pertimbangkan morfisme
% segment-id: o014.li2.chapter2.serre-subcategories-and-k0-groups.g003
$\begin{tikzcd}
	X \arrow[r, "f", shift left] \arrow[r, "g"', shift right] &
	Y \arrow[r, "s \in S"] & W
\end{tikzcd}$,
dengan $sf=sg$. Karena $\Image(f-g)\hookrightarrow\Ker(s)$, maka
$\Image(f-g)\in\Obj(\mathcal{T})$; oleh sebab itu
$Z:=\Ker(f-g)\hookrightarrow X$ merupakan morfisme dalam $S$, yang membuktikan
(S4).

Jadi $S$ merupakan sistem multiplikatif kiri. Semua kondisi bersifat dual,
sehingga $S$ juga merupakan sistem multiplikatif kanan. Sekarang ambil
$Q:\mathcal{A}\to\mathcal{A}/\mathcal{T}:=\mathcal{A}[S^{-1}]$ sebagai
funktor lokalisasi. Proposisi~\ref{prop:localization-extra-property}
menyatakan bahwa $Q$ pada hakikatnya surjektif (bahkan
$\Obj(\mathcal{A})=\Obj(\mathcal{A}/\mathcal{T})$);
Proposisi~\ref{prop:Abel-cat-localization} menyatakan bahwa $Q$ merupakan
funktor eksak di antara kategori-kategori abelian. Perhatikan bahwa
$\identity_{QX}=Q(\identity_X)$ sama dengan $0$ jika dan hanya jika terdapat
$U\in\Obj(\mathcal{A})$ sedemikian sehingga $U\xrightarrow{0}X$ berada dalam
$S$ (ini merupakan latihan sederhana berdasarkan
Korolari~\ref{prop:fraction-zero}); hal ini mengimplikasikan
$X=\Coker[U\xrightarrow{0}X]\in\Obj(\mathcal{T})$. Sebaliknya, jika
$X\in\Obj(\mathcal{T})$, kita dapat mengambil $U=X$. Jadi
$\Ker(Q)=\mathcal{T}$.

Selanjutnya, verifikasi sifat universal. Jika $\Ker(F)\supset\mathcal{T}$,
maka barisan eksak $\Ker(s)\to X\xrightarrow{s\in S}Y\to\Coker(s)$
menunjukkan bahwa $F$ memetakan setiap unsur $S$ menjadi isomorfisme, dan
sifat universal lokalisasi menentukan $G$ yang dicari.

Selanjutnya, verifikasi bahwa $G$ eksak.
Teorema~\ref{prop:localization-additivity} terlebih dahulu memastikan bahwa
$G$ aditif. Karena setiap morfisme di $\mathcal{A}[S^{-1}]$ dapat, setelah
dikomposisikan secara sesuai dengan suatu isomorfisme dari $S$, dipastikan
berasal dari $\mathcal{A}$, maka sifat mempertahankan kernel direduksi menjadi
keeksakan $F$ dan $Q$. Hal yang sama berlaku untuk kokernel. Karena $Q$ pada
hakikatnya surjektif dan $\Ker(Q)=\mathcal{T}$, karakterisasi sifat setia dapat
direduksi menjadi pernyataan (ii) dari Proposisi~\ref{prop:faithful-exact}.
\end{bukti}
% segment-id: o014.li2.chapter2.serre-subcategories-and-k0-groups.p005
Perhatikan bahwa $\mathcal{A}/\mathcal{T}$ dibangun sebagai lokalisasi, sehingga
mungkin merupakan “kategori besar”. Latihan akan memberikan deskripsi lain
tentang $\mathcal{A}/\mathcal{T}$ untuk mengendalikan ukuran kuosien Serre
dengan asumsi $\mathcal{A}$ merupakan kategori well-powered
(Definisi~\ref{def:well-powered}).
\index{masalah ukuran (size issue)}

% segment-id: o014.li2.chapter2.serre-subcategories-and-k0-groups.ex002
\begin{contoh}
Misalkan $S$ merupakan himpunan multiplikatif dari gelanggang komutatif
$R$. Definisikan subkategori penuh $\mathcal{T}$ dari $R\dcate{Mod}$ dengan
$M\in\Obj(\mathcal{T})$ jika dan hanya jika untuk setiap $m\in M$ terdapat
$s\in S$ sedemikian sehingga $sm=0$. Mudah diverifikasi bahwa $\mathcal{T}$
merupakan subkategori Serre. Berikut kita tunjukkan bahwa
$R\dcate{Mod}/\mathcal{T}$ ekuivalen dengan $R[S^{-1}]\dcate{Mod}$.

Pemetaan $M\mapsto M[S^{-1}]:=M\dotimes{R}R[S^{-1}]$ memberikan funktor
eksak $F:R\dcate{Mod}\to R[S^{-1}]\dcate{Mod}$ yang memetakan
$\mathcal{T}$ ke nol, sehingga sifat universal memberikan funktor eksak
$G:R\dcate{Mod}/\mathcal{T}\to R[S^{-1}]\dcate{Mod}$. Jelas bahwa
$\Ker(F)=\mathcal{T}$, sehingga $G$ setia. Selain itu, $G$ pada hakikatnya
surjektif: pandang sebarang modul-$R[S^{-1}]$ $N$ sebagai modul-$R$;
silakan verifikasi isomorfisme modul-$R[S^{-1}]$
$N\dotimes{R}R[S^{-1}]\simeq N$.

Jadi tersisa membuktikan bahwa $G$ penuh dan setia. Diberikan homomorfisme
modul-$R[S^{-1}]$
$\varphi:M_1[S^{-1}]\to M_2[S^{-1}]$, bentuk produk serat
\footnote{Catatan penerjemah (O014-C029): sumber mendefinisikan $M_0$ sebagai
himpunan unsur $m\in M_1$ yang citranya “berasal dari” $M_2$, kemudian
memakai pembatasan $\varphi|_{M_0}:M_0\to M_2$. Prapeta di $M_2$ tidak kanonik
jika $M_2\to M_2[S^{-1}]$ memiliki kernel torsi-$S$, sehingga morfisme tersebut
tidak terdefinisi. Target menggantinya dengan produk serat kanonik beserta
kedua proyeksinya.}
$M_0:=M_1\dtimes{M_2[S^{-1}]}M_2$, di mana pemetaan
$M_1\to M_2[S^{-1}]$ adalah komposisi
$M_1\to M_1[S^{-1}]\xrightarrow{\varphi}M_2[S^{-1}]$.
Tuliskan proyeksi-proyeksinya sebagai $p_1:M_0\to M_1$ dan
$p_2:M_0\to M_2$. Kernel dan kokernel $p_1$ merupakan modul torsi-$S$,
sehingga $p_1$ berada dalam sistem multiplikatif yang menentukan kuosien Serre.
Dengan demikian, atap
% segment-id: o014.li2.chapter2.serre-subcategories-and-k0-groups.g004
$\begin{tikzcd}
	M_1 & M_0 \arrow[l, "p_1"'] \arrow[r, "p_2"] & M_2
\end{tikzcd}$
menentukan morfisme di $R\dcate{Mod}/\mathcal{T}$ yang dipetakan ke $\varphi$.
Ini membuktikan klaim.
\end{contoh}

% segment-id: o014.li2.chapter2.serre-subcategories-and-k0-groups.p006
Sekarang kita memperkenalkan suatu konstruksi penting yang berkaitan erat
dengan subkategori Serre.

% segment-id: o014.li2.chapter2.serre-subcategories-and-k0-groups.d003
\begin{definisi}\label{def:K0}
\index{K0 grup@grup $\mathrm{K}_0$}
\index[sym1]{K0@$\mathrm{K}_0$}
Misalkan $\mathcal{A}$ merupakan kategori abelian. Definisikan grup komutatif
$\mathrm{K}_0(\mathcal{A})$ sebagai berikut, dengan operasi grup ditulis secara
aditif. Pertimbangkan modul-$\Z$ bebas $F$ dengan basis
$\Obj(\mathcal{A})$; tuliskan elemen yang bersesuaian dengan
$X\in\Obj(\mathcal{A})$ sebagai $\lrangle{X}\in F$. Misalkan $R\subset F$
merupakan submodul yang dibangkitkan oleh elemen-elemen berikut:
% segment-id: o014.li2.chapter2.serre-subcategories-and-k0-groups.q002
\[\lrangle{X}-\lrangle{X'}-\lrangle{X''},\quad
0\to X'\to X\to X''\to0:\ \text{barisan eksak pendek di }\mathcal{A},\]
maka $\mathrm{K}_0(\mathcal{A}):=F/R$ disebut \emph{grup $\mathrm{K}_0$}
dari $\mathcal{A}$. Selanjutnya, tuliskan citra
$X\in\Obj(\mathcal{A})$ di $\mathrm{K}_0(\mathcal{A})$ sebagai $[X]$.
\end{definisi}

% segment-id: o014.li2.chapter2.serre-subcategories-and-k0-groups.lm001
\begin{lema}\label{prop:K0-prep}
Untuk setiap kategori abelian $\mathcal{A}$, persamaan-persamaan berikut
berlaku di $\mathrm{K}_0(\mathcal{A})$.
% segment-id: o014.li2.chapter2.serre-subcategories-and-k0-groups.l003
\begin{compactenum}[(i)]
% segment-id: o014.li2.chapter2.serre-subcategories-and-k0-groups.i006
	\item $[0]=0$;
% segment-id: o014.li2.chapter2.serre-subcategories-and-k0-groups.i007
	\item jika $X,Y\in\Obj(\mathcal{A})$ dan $X\simeq Y$, maka $[X]=[Y]$;
% segment-id: o014.li2.chapter2.serre-subcategories-and-k0-groups.i008
	\item $[X\oplus Y]=[X]+[Y]$;
% segment-id: o014.li2.chapter2.serre-subcategories-and-k0-groups.i009
	\item jika
	$0\to X^1\xrightarrow{f^1}\cdots\xrightarrow{f^{n-1}}X^n\to0$
	merupakan barisan eksak di $\mathcal{A}$, maka
	$\sum_{i=1}^{n}(-1)^i[X^i]=0$;
% segment-id: o014.li2.chapter2.serre-subcategories-and-k0-groups.i010
	\item dengan notasi dari \S\ref{sec:Abel-cat-subobjects}, pertimbangkan
	keluarga subobjek $X=X_0\supset\cdots\supset X_r=0$, maka
% segment-id: o014.li2.chapter2.serre-subcategories-and-k0-groups.q003
	\[[X]=\sum_{i=0}^{r-1}[X_i/X_{i+1}].\]
\end{compactenum}
\end{lema}

% segment-id: o014.li2.chapter2.serre-subcategories-and-k0-groups.pf004
\begin{bukti}
Pertimbangkan barisan eksak pendek $0\to0\to0\to0\to0$; ini memberikan (i).
Jika $X\simeq Y$, maka $0\to X\rightiso Y\to0\to0$, bersama dengan (i),
memberikan (ii). Barisan eksak pendek dari
Proposisi~\ref{prop:biproduct-ses} memberikan (iii).

Untuk (iv), tambahkan suku-suku nol dan perpanjang barisan eksak ke kiri dan
kanan hingga tak terhingga. Pertimbangkan barisan eksak pendek
% segment-id: o014.li2.chapter2.serre-subcategories-and-k0-groups.q004
\[0\to\Ker(f^i)\to X^i\to\Image(f^i)\to0,\quad i=1,\ldots,n.\]
Maka
$\left[X^i\right]=\left[\Ker(f^i)\right]+\left[\Image(f^i)\right]
=\left[\Ker(f^i)\right]+\left[\Ker(f^{i+1})\right]$,
lalu ambil jumlah berselang-seling.

Terakhir, (v) diperoleh dengan induksi pada $r$ menggunakan barisan eksak
pendek $0\to X_1\to X_0\to X_0/X_1\to0$.
\end{bukti}

% segment-id: o014.li2.chapter2.serre-subcategories-and-k0-groups.rem001
\begin{catatan}
\index{masalah ukuran (size issue)}
Karena konvensi buku ini adalah merealisasikan grup pada himpunan kecil, dan
mengingat Lema~\ref{prop:K0-prep} (ii), pendekatan yang sepenuhnya ketat adalah
mensyaratkan $\mathcal{A}$ memiliki suatu kerangka kecil
\cite[Lema 2.2.12]{Li1} ketika mendefinisikan
$\mathrm{K}_0(\mathcal{A})$. Hal ini tidak banyak berpengaruh di sini.
\end{catatan}

% segment-id: o014.li2.chapter2.serre-subcategories-and-k0-groups.th002
\begin{teorema}[Prinsip Euler--Poincaré]\label{prop:EP}
Misalkan
$\cdots\to X^i\xrightarrow{d_X^i}X^{i+1}\to\cdots$
merupakan kompleks di kategori abelian $\mathcal{A}$ dengan hanya berhingga
banyak suku tak nol. Maka persamaan
% segment-id: o014.li2.chapter2.serre-subcategories-and-k0-groups.q005
\[\sum_i(-1)^i\left[\Hm^i(X)\right]
=\sum_i(-1)^i\left[X^i\right],\quad
\Hm^i(X):=\Hm\left[X^{i-1}\xrightarrow{d_X^{i-1}}X^i
\xrightarrow{d_X^i}X^{i+1}\right]\]
berlaku di $\mathrm{K}_0(\mathcal{A})$.
\end{teorema}

% segment-id: o014.li2.chapter2.serre-subcategories-and-k0-groups.pf005
\begin{bukti}
Untuk setiap $i$ terdapat barisan eksak
% segment-id: o014.li2.chapter2.serre-subcategories-and-k0-groups.q006
\[0\to\Ker\left(d_X^{i-1}\right)\to X^{i-1}
\xrightarrow{d_X^{i-1}}\Ker\left(d_X^i\right)
\to\Hm^i(X)\to0,\]
dan untuk $|i|\gg0$ semua suku menjadi $0$. Terapkan
Lema~\ref{prop:K0-prep} (iv) dan jumlahkan di
$\mathrm{K}_0(\mathcal{A})$, sehingga diperoleh
% segment-id: o014.li2.chapter2.serre-subcategories-and-k0-groups.q007
\[\sum_i(-1)^i
\left(\left[X^{i-1}\right]+\left[\Hm^i(X)\right]\right)
=\sum_i(-1)^i
\left(\left[\Ker\left(d_X^{i-1}\right)\right]
+\left[\Ker\left(d_X^i\right)\right]\right).\]
Ruas kanan saling meniadakan; dengan menata ulang diperoleh
$\sum_i(-1)^i[\Hm^i(X)]=\sum_i(-1)^i[X^i]$.
\end{bukti}

% segment-id: o014.li2.chapter2.serre-subcategories-and-k0-groups.ex003
\begin{contoh}
Ambil $\mathcal{A}$ sebagai kategori abelian ruang vektor berdimensi hingga
di atas gelanggang pembagian $D$. Karena setiap ruang vektor memiliki basis,
$\mathrm{K}_0(\mathcal{A})=\Z\cdot[D]$. Di sisi lain, pemetaan
$[X]\mapsto\dim_D X$ menentukan homomorfisme grup
$\dim:\mathrm{K}_0(\mathcal{A})\to\Z$ yang memetakan $[D]$ ke $1$.
Jadi $\dim:\mathrm{K}_0(\mathcal{A})\rightiso\Z$.

Jika ruang vektor berdimensi tak hingga diizinkan, grup $\mathrm{K}_0$ akan
menjadi trivial. Alasannya, untuk setiap ruang vektor-$D$ $V$, pertimbangan
kardinalitas himpunan memberikan ruang vektor-$D$ $W$ sedemikian sehingga
$V\oplus W\simeq W$ dan
$\dim_D W=\max\{\dim_D V,\aleph_0\}$; akibatnya $[V]=0$.
\end{contoh}

% segment-id: o014.li2.chapter2.serre-subcategories-and-k0-groups.p007
Untuk funktor eksak $F:\mathcal{A}\to\mathcal{B}$ antara kategori abelian,
karena $F$ mempertahankan barisan eksak pendek, pemetaan $[X]\mapsto[FX]$
menentukan homomorfisme grup
\footnote{Catatan penerjemah (O014-C028): sumber mencetak
$\mathrm{K}_0(f)$, tetapi homomorfisme ini diinduksi oleh funktor $F$;
target menggunakan $\mathrm{K}_0(F)$.}
$\mathrm{K}_0(F):\mathrm{K}_0(\mathcal{A})
\to\mathrm{K}_0(\mathcal{B})$.
Oleh karena itu, jika $\mathcal{A}$ merupakan subkategori abelian dari
$\mathcal{B}$, terdapat homomorfisme alami
$\mathrm{K}_0(\mathcal{A})\to\mathrm{K}_0(\mathcal{B})$.

% segment-id: o014.li2.chapter2.serre-subcategories-and-k0-groups.pr002
\begin{proposisi}\label{prop:K0-exact-seq}
Misalkan $\mathcal{T}$ merupakan subkategori Serre dari kategori abelian
$\mathcal{A}$. Tuliskan kuosien Serre-nya sebagai
$\mathcal{B}:=\mathcal{A}/\mathcal{T}$. Maka funktor-funktor eksak
$\mathcal{T}\to\mathcal{A}\xrightarrow{Q}\mathcal{B}$ menginduksi
homomorfisme yang menghasilkan barisan eksak grup-grup aditif
% segment-id: o014.li2.chapter2.serre-subcategories-and-k0-groups.q008
\[\mathrm{K}_0(\mathcal{T})\to\mathrm{K}_0(\mathcal{A})
\xrightarrow{\mathrm{K}_0(Q)}\mathrm{K}_0(\mathcal{B})\to0.\]
\end{proposisi}

% segment-id: o014.li2.chapter2.serre-subcategories-and-k0-groups.pf006
\begin{bukti}
Pertama, $Q:\mathcal{A}\to\mathcal{B}$ pada hakikatnya surjektif, sehingga
$\mathrm{K}_0(\mathcal{A})\to\mathrm{K}_0(\mathcal{B})$ surjektif.
Selanjutnya, komposisi
$\mathrm{K}_0(\mathcal{T})\to\mathrm{K}_0(\mathcal{A})
\to\mathrm{K}_0(\mathcal{B})$ jelas merupakan $0$. Yang perlu dibuktikan adalah
$\Ker(\mathrm{K}_0(Q))\subset
\Image\left[\mathrm{K}_0(\mathcal{T})\to
\mathrm{K}_0(\mathcal{A})\right]$.

Misalkan $\sum_{i=1}^n a_i[X_i]\in\Ker(\mathrm{K}_0(Q))$, dengan
$a_1,\ldots,a_n\in\Z$. Ini ekuivalen dengan adanya keluarga barisan eksak
pendek di $\mathcal{B}$,
$0\to Y'_j\xrightarrow{f_j}Y_j\xrightarrow{g_j}Y''_j\to0$,
dan $b_j\in\Z$ untuk $j=1,\ldots,m$, sedemikian sehingga persamaan
% segment-id: o014.li2.chapter2.serre-subcategories-and-k0-groups.q009
\begin{equation}\label{eqn:K0-exact-seq-aux-0}
\sum_{i=1}^n a_i\lrangle{QX_i}
=\sum_{j=1}^m b_j
\left(\lrangle{Y_j}-\lrangle{Y'_j}-\lrangle{Y''_j}\right)
\end{equation}
berlaku dalam modul-$\Z$ bebas dengan basis $\Obj(\mathcal{B})$.
Karena $Q:\mathcal{A}\to\mathcal{B}$ dibangun melalui lokalisasi dalam
Teorema~\ref{prop:Serre-quotient}, seperti terlihat dari teorema tersebut
dan \S\ref{sec:cat-localization}, kita mempunyai:
% segment-id: o014.li2.chapter2.serre-subcategories-and-k0-groups.l004
\begin{compactitem}
% segment-id: o014.li2.chapter2.serre-subcategories-and-k0-groups.i011
	\item $Q$ mengidentifikasi $\Obj(\mathcal{A})$ dan
	$\Obj(\mathcal{B})$, sehingga
	\eqref{eqn:K0-exact-seq-aux-0} menghasilkan persamaan dalam
	$\mathrm{K}_0(\mathcal{A})$
% segment-id: o014.li2.chapter2.serre-subcategories-and-k0-groups.q010
	\begin{equation}\label{eqn:K0-exact-seq-aux-1}
	\sum_{i=1}^n a_i[X_i]
	=\sum_{j=1}^m b_j\left([Y_j]-[Y'_j]-[Y''_j]\right);
	\end{equation}
% segment-id: o014.li2.chapter2.serre-subcategories-and-k0-groups.i012
	\item selain itu, setelah menyesuaikan $Y_j$, $Y'_j$, $Y''_j$,
	dan $f_j$, $g_j$ secara tepat dengan isomorfisme-isomorfisme dari
	$S$ di $\mathcal{B}$, kita dapat memastikan adanya morfisme
	$u_j$, $v_j$ di $\mathcal{A}$ sedemikian sehingga
	$f_j=Q(u_j)$ dan $g_j=Q(v_j)$. Berdasarkan
	$Q(v_ju_j)=g_jf_j=0$, dan
	Korolari~\ref{prop:fraction-zero}, kita dapat kembali menyesuaikan
	dengan $S$ sehingga $v_ju_j=0$.
\end{compactitem}

Maka, untuk setiap $1\leq j\leq m$, di $\mathcal{A}$ terdapat barisan eksak
% segment-id: o014.li2.chapter2.serre-subcategories-and-k0-groups.g005
\begin{gather*}
	0\to\Ker(v_j)\to Y_j\xrightarrow{v_j}Y''_j
	\to\Coker(v_j)\to0,\\
	0\to\Ker(u_j)\to Y'_j\xrightarrow{u_j}\Ker(v_j)
	\to\frac{\Ker(v_j)}{\Image(u_j)}\to0.
\end{gather*}
Karena $Q$ eksak, $\Coker(v_j)$, $\Ker(u_j)$, dan
$\frac{\Ker(v_j)}{\Image(u_j)}$ semuanya merupakan objek
$\mathcal{T}=\Ker(Q)$. Oleh karena itu, dalam
$\mathrm{K}_0(\mathcal{A})$ diperoleh
% segment-id: o014.li2.chapter2.serre-subcategories-and-k0-groups.q011
\[[Y_j]-[Y'_j]-[Y''_j]\in
\Image\left[\mathrm{K}_0(\mathcal{T})
\to\mathrm{K}_0(\mathcal{A})\right].\]
Dengan mensubstitusikan persamaan di atas ke
\eqref{eqn:K0-exact-seq-aux-1}, diperoleh
$\Ker(\mathrm{K}_0(Q))\subset
\Image\left[\mathrm{K}_0(\mathcal{T})
\to\mathrm{K}_0(\mathcal{A})\right]$.
\end{bukti}

% segment-id: o014.li2.chapter2.serre-subcategories-and-k0-groups.p008
Menghadapi barisan eksak seperti pada
Proposisi~\ref{prop:K0-exact-seq}, strategi yang selalu berhasil adalah
mencoba memperpanjangnya ke kiri, yakni mencari keluarga grup-$K$ tingkat
tinggi $\mathrm{K}_i(\cdot)$ untuk $i\in\Z_{\geq0}$, beserta barisan eksak
kanonik
% segment-id: o014.li2.chapter2.serre-subcategories-and-k0-groups.q012
\[\cdots\to\mathrm{K}_{i+1}(\mathcal{B})
\to\mathrm{K}_i(\mathcal{T})\to\mathrm{K}_i(\mathcal{A})
\to\mathrm{K}_i(\mathcal{B})\to\cdots\]
Kita juga berharap bahwa keluarga
$(\mathrm{K}_i(\mathcal{A}))_{i\geq0}$ memuat informasi mendalam tentang
$\mathcal{A}$ dan, setidaknya sampai tingkat tertentu, dapat dihitung.
Materi ini termasuk dalam teori K, yang jangkauan penerapannya dapat diperluas
ke kategori eksak, yang lebih umum daripada kategori abelian. Karena konstruksi
terkait didasarkan pada sudut pandang teori homotopi, buku ini tidak dapat
menguraikannya secara rinci; pembaca yang berminat dapat merujuk ke
\cite{Lai19}, atau ringkasan di \cite[\S 13.6]{Bu10}.
\index{kategori eksak (exact category)}
